%%%%%%%%%%%%%%%%%%%%%%%%%%%%%%%%%%%%%%%%%%%%%%%%%%%%%%%%%%%%%%%%%%%%%%
% Lớp 11 | Quan hệ vuông góc trong không gian | Khoảng cách
% Dạng 1H8?5-6: Toán thực tế áp dụng khoảng cách
% Số câu: 2  (sinh tự động bởi tach_cau.py; từ nay thầy sửa trực tiếp file này)
%%%%%%%%%%%%%%%%%%%%%%%%%%%%%%%%%%%%%%%%%%%%%%%%%%%%%%%%%%%%%%%%%%%%%%

%%%%%%%%%%%%%%%% TỰ LUẬN %%%%%%%%%%%%%%%%

% nguồn: Ok_TuanLuan_Toan11_Lan1
\begin{bt}%[1H8V5-6]%[ma:19d1c4f7]
	Một chiếc máy bay cất cánh từ một điểm thuộc mặt đất phẳng nằm ngang. Trong $3$ phút đầu máy  bay bay với vận tốc $500\mathrm{~km}/\mathrm{h}$ và theo đường thẳng tạo với mặt đất một góc $15^\circ$. Hỏi sau $2$ phút, máy bay ở độ cao bao nhiêu kilômét (làm tròn kết quả đến chữ số thập phân thứ nhất)?
	 \loigiai{
	 \immini{
		Sau $2$  phút, máy bay đi được quãng đường là
		$AB=\dfrac{500}{60}\cdot 2=\dfrac{50}{3} (\mathrm{~km})$. \\ 
		Sau $2$ phút,  máy bay ở độ cao là
		 $h=AB\cdot \sin 15^\circ \approx 4, 3(\mathrm{~km})$.
		} 
		{ 
		\begin{tikzpicture}[scale=0.8,line join=round,line cap=round]
				\tkzDefPoints{0/0/M,-1.3/-1.6/N,2.5/-1.6/P} 
				\coordinate (Q) at ($(M)+(P)-(N)$); 
				\coordinate (S) at ($(M)+(0,3)$); 
				\tkzDrawPolygon  (S,N,P,Q)
				\tkzDrawSegments (S,P) 
				\tkzDrawSegments[dashed](M,N M,Q  M,S M,P N,Q)
				\tkzDrawPoints[fill=black,size=4](S,M,N,P,Q) 
				\tkzLabelPoints[above](S) 
				 \tkzLabelPoints[below](N,P)
				 \tkzLabelPoints[right](Q)
				\tkzLabelPoints[left](M) 
			\end{tikzpicture}
		 }
	 }
\end{bt}

% nguồn: Ok_TuanLuan_Toan11_Lan2
\begin{ex}%[1H8V5-6]%[ma:b1ec91fc]
	Một chiếc máy bay cất cánh từ một điểm thuộc mặt đất phẳng nằm ngang. Trong $3$ phút đầu máy bay bay với vận tốc $500\mathrm{~km}/\mathrm{h}$ và theo đường thẳng tạo với mặt đất một góc $15^\circ$. Hỏi sau  $2$ phút, máy bay ở độ cao bao nhiêu kilômét (làm tròn kết quả đến chữ số thập phân thứ nhất)?
	 \loigiai{
	 \immini{
		Sau $2$ phút, máy bay đi được quãng  đường là
		$AB=\dfrac{500}{60}\cdot  2=\dfrac{50}{3} (\mathrm{~km})$. \\
		Sau $2$ phút, máy  bay ở độ cao là
		 $h=AB\cdot \sin 15^\circ \approx 4, 3(\mathrm{~km})$.
		 }
		{ 
		\begin{tikzpicture}[scale=0.8,line join=round,line cap=round]
				 \tkzDefPoints{0/0/M,-1.3/-1.6/N,2.5/-1.6/P}
				\coordinate (Q) at ($(M)+(P)-(N)$); 
				 \coordinate (S) at ($(M)+(0,3)$);
				\tkzDrawPolygon (S,N,P,Q) 
				\tkzDrawSegments  (S,P)
				\tkzDrawSegments[dashed](M,N M,Q  M,S M,P N,Q)
				 \tkzDrawPoints[fill=black,size=4](S,M,N,P,Q)
				\tkzLabelPoints[above](S) 
				\tkzLabelPoints[below](N,P) 
				 \tkzLabelPoints[right](Q)
				\tkzLabelPoints[left](M) 
			\end{tikzpicture}
		} 
	} 
\end{ex}

% nguồn: 2026-11-15_11_1H8_TL2_Vuong_goc (nhập 15/11/2026)
\begin{bt}%[1H8V5-6]%[ma:488643cb]
% nguon: Archive/SP-4Phan-Dot3/Data/KNTT_1-OTGK2_De02.tex #4375

	Kim tự tháp ở Ai Cập có hình dạng là một hình chóp tứ giác đều có đáy là $262$ m và cạnh bên là $230$ m. Giả sử, từ một mặt bên của kim tự tháp ta cần đào một con đường ngắn nhất để đi đến tâm của đáy kim tự tháp, khi đó quãng đường ngắn nhất có độ dài khoảng bao nhiêu mét?
	\loigiai{
		\immini
		{Khoảng cách ngắn nhất từ một điểm tại mặt bên đến tâm $O$ của đáy chính là khoảng cách từ tâm $O$ đến mặt phẳng $(SCD)$.\\
		Dựng $OI\perp CD$ tại $I$, $OH\perp SI$ tại $H$.\\
		Ta có $\heva{&CD\perp OI\\&CD\perp SO} \Rightarrow CD\perp (SOI) \Rightarrow OI\perp OH$.\\
		Ta có $\heva{& OH\perp SI\\& OH\perp CD} \Rightarrow OH\perp (SCD)$.}
		{\begin{tikzpicture}[declare function={a=2.65;b=4.5;h=4;}, line join=round, line cap=round, >=stealth,scale=0.9]
			\path (0,0) coordinate (A)
			(-135:a) coordinate (B)
			(b,0) coordinate (D)
			($(D)-(A)+(B)$) coordinate (C)
			($(B)!0.5!(D)$) coordinate (O)
			($(O)+(0,h) $) coordinate (S)
			($(C)!.5!(D)$) coordinate (I)
			($(S)!0.65!(I)$) coordinate (H);		
			\path ($(D)-(A)+(B)$) coordinate (C);
			\draw[dashed] (S)--(A)--(B) (A)--(D) (A)--(C) (S)--(O) (B)--(D) (O)--(I) (O)--(H) ;
			\draw (B)--(C)--(D) (B)--(S) (D)--(S)--(C) (S)--(I);
			\draw pic[draw=black,angle radius=0.2cm] {right angle = O--I--C};
			\draw pic[draw=black,angle radius=0.2cm] {right angle = S--H--O};
			\foreach \t/\g in {A/150,B/-90,C/-90,D/0,S/90,O/-90,I/-20,H/10}{\draw[fill=black] (\t) circle (0.75pt) node[shift={(\g:7pt)}]{$\t$};}
		\end{tikzpicture}}
		\noindent
		Vậy $OH = \mathrm{d}\left(O,(SCD)\right)$.\\
		Ta có $OC=\dfrac{262\sqrt{2}}{2} \Rightarrow SO=\sqrt{SC^2-OC^2}=\sqrt{18578}$.\\
		Ta có $OH=\dfrac{OI\cdot SO}{\sqrt{OI^2+SO^2}} = \dfrac{131\cdot\sqrt{18578}}{\sqrt{131^2+18578}}$ (mét).
	}
\end{bt}
